\chapter{Background and Related Work}
\label{ch:background}

This chapter establishes the technical and historical context against which the project's choices and findings should be read. \Cref{sec:bg-history} surveys the development of terrain classification for mobile robots, with emphasis on proprioceptive sensing as an alternative to camera- and LiDAR-based perception, and \Cref{sec:bg-proprio} examines the physical rationale for using inertial and wheel-encoder signals. \Cref{sec:bg-features-ml} then reviews the windowing, feature-engineering, and supervised-learning techniques applied to short segments of vibration data, covering both classical hand-crafted pipelines and end-to-end deep networks. \Cref{sec:bg-vel} discusses the velocity--vibration confound that motivates the explicit speed-regime analysis carried out later in the thesis, and \Cref{sec:bg-transition} introduces statistical change-point detection on classifier posteriors, which underpins the CUSUM-based transition detector implemented in this work. Finally, \Cref{sec:bg-agri} relates the project to agricultural field robotics, and \Cref{sec:bg-summary} summarises the gap that this thesis addresses.

\section{Terrain Classification for Mobile Robots}
\label{sec:bg-history}

Mobile robots operating outside structured environments must continuously adapt to the surface they are driving on: loose soil and compacted gravel impose different rolling resistances, demand different control gains, and pose very different risks of slip or immobilisation. Proprioceptive terrain recognition therefore sits at the boundary between perception and control, and has been an active research topic since the early 2000s.

The first widely cited proprioceptive approach is the work of Brooks and Iagnemma~\cite{Brooks2005}, who showed that the short-time vibration spectrum recorded by an IMU on a wheeled platform encodes enough information to discriminate gravel, grass, sand and concrete in real time from one-second windows of vertical acceleration. The idea was subsequently extended through frequency-response models of the wheel--terrain contact~\cite{DuPont2008} and applied in planetary robotics, where dust and shadowing limit the reliability of cameras and LiDAR~\cite{Weiss2011,Otsu2016}. A similar emphasis on intrinsic sensing appears in slip-detection on wheeled rovers~\cite{Ojeda2006} and supervised classification on legged platforms~\cite{Bermudez2012}.

Across this body of work the basic formulation has been remarkably stable: slice the signal into short windows, summarise each window with a fixed-size feature vector, and train a supervised classifier. The sensor set (encoders, motor current, suspension travel, microphones) and the classifier family ($k$-NN, SVMs, random forests, gradient boosting, neural networks) have grown, but the template has not. The pipeline developed in this thesis (\Cref{ch:processing,ch:feature,ch:classification}) follows the same template, adapted to the characteristics of a small, lightweight agricultural robot.



\section{Proprioceptive Sensing for Terrain Recognition}
\label{sec:bg-proprio}

Robotics literature distinguishes \emph{exteroceptive} sensing (cameras, LiDAR, radar, GNSS), which observes the environment from a distance, from \emph{proprioceptive} or \emph{intrinsic} sensing (IMU, wheel encoders, motor current, joint torques), which observes the robot's own response to that environment. The same wheel on grass versus concrete produces detectably different IMU traces because the contact mechanics differ in stiffness, damping and contact-patch geometry~\cite{Iagnemma2004,Brooks2005}.

Several practical advantages motivate the intrinsic approach over vision-based alternatives:
\begin{itemize}
    \item It is invariant to lighting, fog, dust, glare and visual occlusion, all of which routinely defeat camera-based methods in outdoor and agricultural environments~\cite{Reina2016}.
    \item It directly senses the property of interest, the dynamic interaction between the wheel and the surface, rather than a visual proxy such as colour or texture.
    \item It requires no additional sensors beyond those already needed for inertial odometry and motor control, keeping the cost and power budget low.
\end{itemize}

The disadvantages are equally well documented. Proprioceptive signals are inherently local: they describe only the patch of ground currently in contact with the wheels, with no look-ahead capability. They are also strongly coupled to robot kinematics: speed, heading, payload and suspension state all modulate the same vibration that we seek to attribute to the surface~\cite{Bermudez2012,Coyle2008}. \Cref{sec:bg-vel} returns to this confound, which is treated explicitly through an explicit speed-regime analysis on the campus dataset.

A typical IMU reports three-axis linear acceleration and three-axis angular rate at \SI{100}{\hertz} and above. On a wheeled robot the vertical acceleration channel carries most of the terrain signature, because surface irregularities couple primarily into vertical motion through the suspension. The lateral and longitudinal channels add directional information about side-slip and traction events. Body-frame pitch and roll rates capture the pitching and rocking induced by uneven contact across the wheelbase. Yaw rate, by contrast, is dominated by steering and is therefore conventionally excluded from terrain-recognition features~\cite{Brooks2005}; this thesis follows that convention, as detailed in \Cref{ch:feature}.

\section{Time-Series Classification: Features and Models}
\label{sec:bg-features-ml}

Raw IMU traces are non-stationary, sensor-specific and high-dimensional, so classifiers are rarely trained directly on the unprocessed samples. The classical pipeline first \emph{windows} the signal into short, overlapping segments, typically between \SI{0.5}{\second} and \SI{4}{\second} in length~\cite{Brooks2005,Weiss2011}, and then summarises each window with a fixed set of features that are approximately stationary within the window. Three families dominate the literature:

\begin{itemize}
    \item \textbf{Time-domain statistics} include mean, standard deviation, root-mean-square (RMS), peak, range, interquartile range, skewness, kurtosis and zero-crossing rate; Hjorth~\cite{Hjorth1970} added the \emph{activity}, \emph{mobility} and \emph{complexity} parameters, which provide compact estimates of signal power, mean frequency and bandwidth without an explicit Fourier transform. 
    
    Concretely, for a window $x[n]$ of length $N$, let $\dot{x}[n] = x[n]-x[n-1]$ and $\ddot{x}[n] = \dot{x}[n]-\dot{x}[n-1]$ denote the first and second discrete differences. The Hjorth parameters are then defined as:
    \begin{align} A(x) &= \mathrm{var}(x), \label{eq:hjorth-activity} \\ M(x) &= \sqrt{\frac{\mathrm{var}(\dot{x})}{\mathrm{var}(x)}}, \label{eq:hjorth-mobility} \\ C(x) &= \frac{M(\dot{x})}{M(x)} = \sqrt{\frac{\mathrm{var}(\ddot{x})\,\mathrm{var}(x)}{\mathrm{var}(\dot{x})^{2}}}, \label{eq:hjorth-complexity}
    \end{align}
    representing signal power (activity), an estimate of mean frequency (mobility), and an estimate of bandwidth (complexity).

    \item \textbf{Frequency-domain features} include magnitudes of the short-time Fourier transform, power in physically motivated frequency bands, typically \SIrange{1}{5}{\hertz} for low-frequency body modes, \SIrange{5}{20}{\hertz} for terrain texture, and \SIrange{20}{40}{\hertz} for high-frequency wheel--surface interaction, together with spectral centroid and spectral entropy~\cite{Welch1967,Brooks2005}; jerk RMS, the root-mean-square of the first time derivative of acceleration, captures abrupt impacts and is highly discriminative between hard and compliant surfaces. 

    Letting $\hat{X}(f)$ denote the discrete Fourier transform of the windowed signal and $S_{xx}(f) = |\hat{X}(f)|^{2}/N$ its periodogram estimate of the one-sided power spectral density~\cite{Welch1967}, the power in a band $[f_{\mathrm{lo}},\,f_{\mathrm{hi}}]$ is
    \begin{equation}
        P_{[f_{\mathrm{lo}},\,f_{\mathrm{hi}}]}(x)
          \;=\; \int_{f_{\mathrm{lo}}}^{f_{\mathrm{hi}}} S_{xx}(f)\,\mathrm{d}f.
        \label{eq:band-power}
    \end{equation}
    The spectral centroid $f_{\mathrm{c}}$ and the (normalised) spectral entropy $H_{\mathrm{s}}$ are computed from the same periodogram as
    \begin{equation}
        f_{\mathrm{c}}(x) = \frac{\sum_{f>0} f\, S_{xx}(f)}{\sum_{f>0} S_{xx}(f)},
        \qquad
        H_{\mathrm{s}}(x) = -\sum_{f>0} p(f)\log_{2} p(f),
        \quad p(f) = \frac{S_{xx}(f)}{\sum_{f'>0} S_{xx}(f')}.
        \label{eq:spectral}
    \end{equation}
    Jerk RMS, finally, captures the rate of change of acceleration; for a sampling period $T_{s}$ it is computed as
    \begin{equation}
        \mathrm{JerkRMS}(x)
          \;=\; \sqrt{\,\frac{1}{N-1}\sum_{n=1}^{N-1}\!\left(\frac{x[n]-x[n-1]}{T_{s}}\right)^{2}}.
        \label{eq:jerk-rms}
    \end{equation}

    \item \textbf{Cross-axis and cross-sensor correlations}, finally, encode the joint linear--rotational signature of, for instance, a wheel sinking into mud or cresting a bump.

\end{itemize}

The feature set grows quickly: a dozen descriptors per axis, six IMU axes plus magnitude channels, several frequency bands and a handful of correlations can easily exceed one hundred features. The standard recipe, also used in this thesis, is to rank features by mutual information against the class label~\cite{Vergara2014,Cover2006,Bao2004} and then remove the worst redundancies by Pearson correlation pruning. Mutual information measures the reduction in uncertainty about the class label $Y$ given a feature $X$,
\begin{equation}
    I(X;Y) \;=\; \sum_{y \in \mathcal{Y}} \int_{\mathcal{X}}
                p(x,y)\,\log\frac{p(x,y)}{p(x)\,p(y)}\,\mathrm{d}x
          \;=\; H(Y) - H(Y \mid X),
    \label{eq:mi}
\end{equation}
where $H(\cdot)$ is Shannon entropy~\cite{Cover2006}; a feature with high MI is informative about the terrain label irrespective of any specific classifier, which makes it a model-agnostic ranking criterion.

Once a window has been reduced to a feature vector, classification becomes a generic supervised learning problem. The four classical baselines used in this thesis --- logistic regression~\cite{Hosmer2013}, random forest~\cite{Breiman2001}, gradient-boosted trees (XGBoost)~\cite{Chen2016} and a support vector machine with an RBF kernel~\cite{Cortes1995} --- span the standard linear, ensemble and kernel-method axes; each places different priors on the data (globally linear class boundaries; threshold-based non-linear feature interactions; local similarity in feature space).

Deep neural networks have transformed many time-series classification tasks since the mid-2010s~\cite{LeCun2015,Fawaz2019}, and for human-activity recognition, the closest cousin of terrain classification, one-dimensional convolutional networks (1D-CNNs) on the raw multi-channel IMU stream routinely match or exceed hand-crafted feature pipelines~\cite{Wang2017,Fawaz2019}. Recurrent architectures (LSTM~\cite{Hochreiter1997}, GRU) and more recent attention-based variants extend the receptive field across windows but generally demand larger training datasets to regularise effectively. Two findings recur for terrain classification specifically: classical methods with well-designed features remain competitive at the moderate dataset sizes typical of field robotics (tens of thousands of windows rather than millions)~\cite{Weiss2011,Otsu2016,Fawaz2019}, and the deep models gain the most when the input is rich and high-dimensional (multiple IMUs, audio or vision streams). These observations motivate the comparison used here: the four classical baselines are evaluated against an MLP on the engineered features and a 1D-CNN on raw windows of the five terrain-discriminative IMU channels, under identical Leave-One-Run-Out cross-validation (\Cref{ch:classification}).

\section{Velocity Dependence and Motion Confounds}
\label{sec:bg-vel}

A recurring difficulty in proprioceptive terrain classification and the reason for the dataset split adopted in this thesis is that the same features which respond to terrain also respond to robot velocity, payload and suspension dynamics~\cite{Bermudez2012,Coyle2008}. Driving faster over a given surface produces a higher-amplitude, broader-spectrum vibration signature; a heavier payload damps high-frequency content; a stiff versus compliant suspension shifts the resonance band. A classifier trained on logs that combine smooth terrain at low speed with rough terrain at high speed can therefore learn to discriminate the classes from speed alone, leaving the terrain information unused, and reported test accuracies in such studies are correspondingly optimistic.

Mitigation strategies fall into three groups: (i) \emph{conditioning on speed} by including wheel-velocity statistics as additional features and letting the model learn the joint dependence~\cite{Ojeda2006}; (ii) \emph{normalising the features} by dividing band-power features by a velocity-dependent baseline so that residual variation can be attributed to terrain; and (iii) \emph{restricting the speed range} to a narrow, physically meaningful band before training and evaluation~\cite{Bermudez2012}.

This thesis uses approaches (i) and (iii), but for distinct reasons on the two datasets. On the \emph{farm}, the robot operates at an intrinsically low and fairly constant speed: its task is weeding and managing intruding plants between crop rows, which is mechanically incompatible with fast traversal. A low-speed, near-constant-velocity regime is therefore the natural deployment condition, not an artificial restriction, and the farm experiments are conducted in that regime. On the \emph{campus} development route, the robot was driven at low speed within each terrain segment, simulating the farm operating conditions for which the pipeline was being developed, and at higher speed when moving between segments. The resulting mix produced enough windows in both regimes that a full-speed and a low-speed subset can be evaluated side by side, which is useful for quantifying how much of the apparent classification accuracy is driven by speed-related vibration rather than by terrain identity. The results of both comparisons are reported in \Cref{ch:results}.

\section{Terrain Transition and Change-Point Detection}
\label{sec:bg-transition}

A classifier that operates on a sliding window provides, at each step, a posterior distribution over the terrain classes. The simplest \emph{implicit} change-point detector takes the hard argmax of this posterior and declares a transition whenever the predicted label changes; with a smoothing rule that requires $n_{\text{stable}}$ consecutive consistent windows the detector acts as a low-pass filter on the prediction stream. This approach is widely used in practice because it requires no additional machinery, but its detection delay grows linearly with $n_{\text{stable}}$, its false-alarm rate is hard to tune analytically, and it discards the soft information in the posterior.

Statistical change-point detection offers a principled alternative. The cumulative-sum (CUSUM) test of Page~\cite{Page1954} accumulates a log-likelihood ratio between two candidate distributions and raises an alarm when the running sum exceeds a fixed threshold; classical analysis ties the threshold to a target average run length under the null hypothesis. Blanke et al.~\cite[\S6.4]{Blanke2003} present the bank-of-CUSUMs formulation used in the fault-detection-and-isolation literature, in which one CUSUM is maintained for each candidate post-change hypothesis and the alarm class is the argmax of the running statistics, exactly the structure required for multiclass terrain transitions. Bayesian online change-point detection (BOCD)~\cite{Adams2007} is a fully probabilistic alternative, and the survey of Truong et al.~\cite{Truong2020} reviews these and related approaches in a unified framework.

Concretely, let $\mathbf{z}(k) \in \Delta^{C-1}$ denote the classifier's posterior over $C$ terrain classes at window index $k$, and let $c_{0}$ be the currently active class (initialised from the first stationary windows of a run and updated after each alarm). At each time $k$ the detector tests, for every candidate post-change class $c \neq c_{0}$, the hypotheses
\begin{align}
  \mathcal{H}_0       &:\; \mathbf{z}(i) \sim p_{c_0}
                          \quad \text{for } 1 \le i \le k, \\
  \mathcal{H}_1^{(c)} &:\; \mathbf{z}(i) \sim p_{c_0}
                          \text{ for } i < k_0, \;
                          \mathbf{z}(i) \sim p_{c}
                          \text{ for } i \ge k_0,
\end{align}
where the change time $k_0$ is unknown (Problem~6.6 of \cite{Blanke2003}). The per-step log-likelihood ratio between these hypotheses is
\begin{equation}
    s_{c}(k) \;=\; \ln \frac{p_{c}(\mathbf{z}(k))}{p_{c_{0}}(\mathbf{z}(k))},
    \label{eq:cusum-llr}
\end{equation} and the recursive CUSUM accumulates positive evidence while discarding negative drift~\cite[\S6.4.2]{Blanke2003}:
\begin{equation}
    g_{c}(k) \;=\; \max\!\bigl(0,\;g_{c}(k-1) + s_{c}(k)\bigr),
    \qquad g_{c}(0) = 0.
    \label{eq:cusum-recursion}
\end{equation}
An alarm is raised, and a transition to class $\hat{c}$ declared, the first time any bank statistic crosses the threshold $h$:
\begin{equation}
    k_{a} \;=\; \min\Bigl\{ k \,:\, \max_{c \neq c_{0}} g_{c}(k) > h \Bigr\},
    \qquad
    \hat{c} \;=\; \argmax_{c \neq c_{0}} g_{c}(k_{a}).
    \label{eq:cusum-alarm}
\end{equation}
The estimated change time is recovered as the index at which the cumulative sum $S_{\hat{c}}(k) = \sum_{i=1}^{k} s_{\hat{c}}(i)$ last reached its minimum before the alarm; after every alarm the active class is updated to $c_0 \leftarrow \hat{c}$ and all bank statistics are reset to zero (Algorithm~6.3 of \cite{Blanke2003}).

The threshold $h$ controls the trade-off between detection delay and false-alarm rate. Wald's approximation~\cite[Eq.~6.74]{Blanke2003} shows that the mean time between false alarms (MTBFA) grows approximately as $\exp(h)$ under the null hypothesis while the mean detection delay grows approximately linearly in $h$. This lets $h$ be calibrated on stationary segments of the training data to a chosen MTBFA and then used unchanged at evaluation, an analytic handle that the implicit $n_{\text{stable}}$ detector lacks. An alternative formulation independent of the classifier applies the same recursion to the standardised IMU feature vector under a Gaussian class-hypothesis model~\cite[\S6.4.3]{Blanke2003}, and is used in \Cref{ch:results} as a baseline detector whose performance does not depend on classifier accuracy.

In the proprioceptive terrain-classification literature, transition detection is most often left implicit, with only a handful of studies addressing the problem explicitly~\cite{Bai2019}. This thesis adopts the bank-of-CUSUMs formulation on the classifier's soft posterior stream as the principal explicit detector, and compares it with both the implicit $n_{\text{stable}}$ detector and the vector-CUSUM baseline in \Cref{ch:results}.


\section{Agricultural Field Robotics}
\label{sec:bg-agri}

The application context of this thesis is the DTU Robotic Intercropping project~\cite{RoboticIntercroppingDTU2026,RoboticIntercroppingProject2026}, in which a small wheeled platform operates among row crops in an agricultural field. Agricultural field robotics differs from the urban-driving and planetary-rover settings that dominate the terrain-classification literature in two respects that motivate the present work.

First, the relevant surfaces --- grass between rows, tilled soil within the plot, and compacted dirt access tracks --- can be hard to discriminate visually under field conditions: variable outdoor lighting, shadowing, and the dust and debris that accumulate on or are flicked onto camera lenses during driving all degrade the reliability of vision-based methods~\cite{Reina2016}. Discriminating these surfaces nevertheless matters for traction control and crop-safe manoeuvring, which makes the proprioceptive vibration signature a useful complementary channel.

Second, mainstream agricultural-robotics research has focused on exteroceptive sensing such as stereo vision, RGB-D and LiDAR for tasks such as row following and crop-versus-weed discrimination~\cite{Reina2016}. Proprioceptive terrain recognition in the intercropping context has received comparatively little attention. This thesis aims to fill part of that gap by providing a labelled, multi-run farm dataset (\Cref{ch:collection}) and a benchmarked pipeline tailored to the agricultural setting.

\section{Position of this Thesis}
\label{sec:bg-summary}

The literature surveyed above establishes that (i) IMU-based proprioceptive terrain classification is a mature technique with strong empirical results on outdoor and planetary surfaces; (ii) carefully engineered features remain competitive with deep models at moderate dataset sizes; (iii) the speed--vibration confound is a known but rarely isolated source of optimistically biased accuracy estimates; and (iv) explicit transition detection on the classifier's posterior stream is principled but under-used in this application area.

Against that backdrop, this thesis makes three contributions: a labelled multi-run dataset from a wheeled agricultural robot covering three farm and five campus terrain classes; a controlled comparison of six classifier families --- logistic regression, random forest, gradient-boosted trees, support vector machine, multilayer perceptron and 1D-CNN --- on both full-speed and low-speed dataset variants, isolating the speed contribution to apparent accuracy; and an online bank-of-CUSUMs transition detector on the soft posterior stream, evaluated alongside the implicit sliding-window detector. The remainder of the report develops each of these choices in turn (\Cref{ch:setup,ch:collection,ch:processing,ch:feature,ch:classification,ch:results}).